% ============================================================================
%  PÓSTER — Modelado numérico de la histéresis aerodinámica (NACA 0018)
%  CC2104 Métodos Numéricos — UTEC — 2026-2 — Grupo 1
%
%  Formato: 48 × 36 pulgadas (horizontal), pensado para impresión.
%  Compilación: pdfLaTeX, dos pasadas (Palatino + Helvetica).
%  Comando: pdflatex poster_NACA0018.tex  (repetir para estabilizar el overlay)
%
%  Las figuras 2–5 todavía no están incluidas: en la columna 3 hay espacios
%  reservados con el tamaño y la ubicación definitivos. Para reemplazarlos,
%  sigue las instrucciones del archivo Figuras/LEEME.md.
% ============================================================================

\documentclass[final]{article}

\usepackage[paperwidth=121.92cm, paperheight=91.44cm,
            top=13.4cm, bottom=3.4cm, left=2.4cm, right=2.4cm]{geometry}

\usepackage[utf8]{inputenc}
\usepackage[T1]{fontenc}
\usepackage[spanish,es-noshorthands]{babel}
\usepackage{amsmath}
\usepackage{mathpazo}                 % cuerpo serif (Palatino)
\usepackage[scaled=0.92]{helvet}      % títulos sans (Helvetica)
\usepackage{textcomp}

\usepackage{graphicx}
\usepackage{tikz}
\usetikzlibrary{positioning,arrows.meta,calc}
\usepackage{booktabs}
\usepackage{array}
\usepackage{colortbl}
\usepackage{enumitem}
\usepackage{xcolor}
\usepackage{microtype}
\usepackage[hidelinks]{hyperref}

\pagestyle{empty}
\setlength{\parindent}{0pt}
\setlength{\parskip}{0.24cm}
\setlength{\abovedisplayskip}{0.34cm}
\setlength{\belowdisplayskip}{0.34cm}
\setlength{\abovedisplayshortskip}{0.20cm}
\setlength{\belowdisplayshortskip}{0.20cm}

% ----------------------------------------------------------------------------
% Tipografías
% ----------------------------------------------------------------------------
% Cuerpo: Palatino (mathpazo). Títulos: Helvetica.
\newcommand{\headingfont}{\sffamily}
\newcommand{\headingfontlight}{\sffamily}

% ----------------------------------------------------------------------------
% Paleta UTEC
% ----------------------------------------------------------------------------
\definecolor{utecNavy}{HTML}{00274C}
\definecolor{utecCyan}{HTML}{19B5E0}
\definecolor{utecGray}{HTML}{3A3A3A}
\definecolor{utecLight}{HTML}{EAF3F8}

% ----------------------------------------------------------------------------
% Escalas tipográficas
% ----------------------------------------------------------------------------
\newcommand{\cuerpo}{\fontsize{22.5}{28.5}\selectfont}
\newcommand{\cuerpoSmall}{\fontsize{19}{24}\selectfont}
\newcommand{\captionfont}{\fontsize{14.5}{18}\selectfont}

% ----------------------------------------------------------------------------
% Encabezado de bloque (título centrado + regla, como el póster de referencia)
% ----------------------------------------------------------------------------
\newcommand{\bloque}[1]{%
  \par\vspace{0.10cm}%
  \begin{center}
    {\headingfont\bfseries\color{utecNavy}\fontsize{29}{34}\selectfont #1}
  \end{center}%
  \vspace{-0.44cm}%
  {\color{utecNavy}\rule{\linewidth}{1.1pt}}\par
  \vspace{0.34cm}%
}

\newcommand{\subbloque}[1]{%
  \par\vspace{0.26cm}%
  {\headingfont\bfseries\color{utecNavy}\fontsize{21.5}{25}\selectfont #1}\par
  \vspace{0.10cm}%
}

% ----------------------------------------------------------------------------
% Bloque de figura con imagen real
% #1 = archivo, #2 = número de figura, #3 = pie, #4 = altura máxima
% ----------------------------------------------------------------------------
\newcommand{\figimage}[4]{%
  \begin{minipage}[t]{18.05cm}
    \centering
    \includegraphics[width=\linewidth,height=#4,keepaspectratio]{#1}\\[0.24cm]
    {\captionfont\textbf{Figura #2.} #3\par}
  \end{minipage}%
}

% ----------------------------------------------------------------------------
% Estilos de tabla
% ----------------------------------------------------------------------------
\newcommand{\tabhead}[1]{\headingfont\bfseries\color{utecNavy}#1}
\renewcommand{\arraystretch}{1.18}

% ============================================================================
\begin{document}

% ============================================================================
% BANDA SUPERIOR
% ============================================================================
\begin{tikzpicture}[remember picture, overlay]
  \fill[utecNavy] (current page.north west) rectangle
    ([yshift=-12cm]current page.north east);

  \node[anchor=north, text width=100cm, align=center,
        text=white, font=\headingfont\bfseries, yshift=-1.85cm]
    at (current page.north) {\fontsize{56}{64}\selectfont
      \spaceskip=0.28em
      Modelado numérico de la histéresis aerodinámica asociada a la pérdida};

  \node[anchor=north, text width=100cm, align=center,
        text=white, font=\headingfont\bfseries, yshift=-4.22cm]
    at (current page.north) {\fontsize{56}{64}\selectfont
      \spaceskip=0.28em
      de un ala NACA 0018 mediante integración y continuación};

  \node[anchor=north, text width=108cm, align=center,
        text=white, font=\headingfont, yshift=-7.30cm]
    at (current page.north) {\fontsize{25}{31}\selectfont
      Jesús A. Salirrosas Vásquez \textperiodcentered\ Juan Diego Mejía Armas
      \textperiodcentered\ Jonathan L. Reyes Torres \textperiodcentered\
      Miguel A. Ramos Bonilla \textperiodcentered\ Moises Suni Chino};

  \node[anchor=north, text width=108cm, align=center,
        text=white, font=\headingfontlight, yshift=-9.05cm]
    at (current page.north) {\fontsize{21}{25}\selectfont
      Universidad de Ingeniería y Tecnología (UTEC)
      \textperiodcentered\ CC2104 Métodos Numéricos};

  % Logo UTEC (archivo extraído del informe; se puede reemplazar por el oficial)
  \node[anchor=north east, xshift=-1.8cm, yshift=-1.55cm,
        fill=white, rounded corners=0.18cm, inner sep=0.28cm]
    at (current page.north east) {\includegraphics[height=7.2cm]{Figuras/logo_utec.png}};

  % Línea de acento cian
  \fill[utecCyan] ([yshift=-11.82cm]current page.north west) rectangle
    ([yshift=-12cm]current page.north east);
\end{tikzpicture}

% ============================================================================
% PIE DE PÁGINA
% ============================================================================
\begin{tikzpicture}[remember picture, overlay]
  \fill[utecNavy] (current page.south west) rectangle
    ([yshift=2.2cm]current page.south east);

  \node[anchor=west, text=white, font=\headingfont,
        xshift=2.4cm, yshift=1.10cm]
    at (current page.south west) {\fontsize{17}{17}\selectfont 2026-2};

  \node[anchor=center, text=white, font=\headingfont,
        xshift=0cm, yshift=1.10cm]
    at (current page.south) {\fontsize{17}{17}\selectfont
      CC2104 \textperiodcentered\ Métodos Numéricos};

  \node[anchor=east, text=white, font=\headingfont,
        xshift=-2.4cm, yshift=1.10cm]
    at (current page.south east) {\fontsize{17}{17}\selectfont
      ABP \textperiodcentered\ Grupo 1};
\end{tikzpicture}

% ============================================================================
% CUERPO — TRES COLUMNAS
% ============================================================================
\noindent
% ============================================================================
% COLUMNA 1
% ============================================================================
\begin{minipage}[t]{37.44cm}
\vspace{0pt}
\cuerpo
\raggedright

\bloque{Introducción}

La \textbf{pérdida aerodinámica} ocurre cuando un perfil alar supera
determinados ángulos de ataque: el flujo se desprende y las cargas
aerodinámicas disminuyen. La \textbf{histéresis aerodinámica} aparece cuando,
para un mismo ángulo de ataque, la sustentación depende de si el ángulo está
aumentando o disminuyendo. El sistema presenta entonces dos ramas estables y
una transición entre ellas.

En este trabajo se estudia numéricamente la histéresis asociada a la pérdida
de un ala con perfil \textbf{NACA 0018}, mediante un modelo no lineal que
representa la adherencia y el desprendimiento del flujo.

\bloque{Objetivos}

\textbf{Objetivo general.} Modelar numéricamente la histéresis aerodinámica
asociada a la pérdida de un ala NACA 0018, analizando la transición entre las
ramas estables de las cargas aerodinámicas.

\vspace{0.12cm}
\textbf{Objetivos específicos}
\begin{itemize}[leftmargin=0.80cm, itemsep=0.16cm, topsep=0.12cm,
                  label=\textcolor{utecCyan}{\textbullet}]
  \item Implementar un modelo no lineal de relajación--retardo para la
        dinámica de adherencia del flujo.
  \item Evaluar el comportamiento del modelo mediante integración numérica.
  \item Analizar la sensibilidad ante variaciones de frecuencia, amplitud y
        parámetros temporales $b$ y $\tau_2$.
  \item Estimar el coeficiente de fuerza normal $C_N$ en condiciones estáticas
        y dinámicas, contrastándolo con datos experimentales digitalizados.
\end{itemize}

\vspace{0.30cm}
\begin{tikzpicture}
  \node[draw=utecCyan, fill=utecLight, line width=1.2pt,
        rounded corners=0.20cm, text width=33.6cm, inner sep=0.42cm,
        align=left]
    {\headingfont\bfseries\color{utecNavy}\fontsize{18}{22}\selectfont
       Pregunta de investigación\\[0.14cm]
     {\rmfamily\mdseries\fontsize{19.5}{24.5}\selectfont
       ¿Cómo se manifiesta la histéresis aerodinámica de un ala NACA 0018 y
       cómo cambia la transición entre sus ramas estables al variar la
       frecuencia y los parámetros temporales del modelo?}};
\end{tikzpicture}

\bloque{Metodología}

\begin{center}
\begin{tikzpicture}[
  box/.style={draw=utecNavy, fill=utecLight, rounded corners=0.18cm,
              line width=1pt, text width=29.5cm, align=center,
              minimum height=2.28cm, inner sep=0.18cm,
              font=\fontsize{16.5}{19}\selectfont},
  small/.style={draw=utecNavy, fill=utecLight, rounded corners=0.15cm,
                line width=1pt, text width=8.4cm, align=center,
                minimum height=1.65cm, inner sep=0.12cm,
                font=\fontsize{15.5}{18}\selectfont},
  arrow/.style={-{Latex[length=3mm,width=2.6mm]}, line width=1.4pt,
                draw=utecNavy},
  line/.style={line width=1.4pt, draw=utecNavy},
  node distance=0.90cm
]

\node[box] (n1) {Inicio};
\node[box, below=of n1] (n2) {Digitalización de datos experimentales\\
  \emph{44 puntos estáticos \textperiodcentered\ 31 puntos dinámicos}};
\node[box, below=of n2] (n3) {Ajuste de parámetros\\
  \emph{búsqueda global + mínimos cuadrados}};
\node[box, below=of n3] (n4) {Formulación del modelo diferencial no lineal};
\node[box, below=of n4] (n5) {Análisis estático\\
  \emph{Newton--Raphson + continuación pseudo-longitud de arco}};
\node[box, below=of n5] (n6) {Análisis dinámico: integración temporal};

\node[small, below=1.35cm of n6, xshift=-9.6cm] (e) {Euler};
\node[small, below=1.35cm of n6]               (h) {Heun};
\node[small, below=1.35cm of n6, xshift=9.6cm]  (r) {RK4};

\node[box, below=1.35cm of h] (n7) {Análisis de sensibilidad\\
  \emph{frecuencia, amplitud, $b$ y $\tau_2$}};
\node[box, below=of n7] (n8) {Contraste con datos digitalizados y cálculo del RMSE};
\node[box, below=of n8] (n9) {Conclusiones};

\draw[arrow] (n1) -- (n2);
\draw[arrow] (n2) -- (n3);
\draw[arrow] (n3) -- (n4);
\draw[arrow] (n4) -- (n5);
\draw[arrow] (n5) -- (n6);

\coordinate (split) at ($(n6.south)+(0,-0.90)$);
\draw[line] (n6.south) -- (split);
\draw[line] (split -| e.north) -- (split -| r.north);
\draw[arrow] (split -| e.north) -- (e.north);
\draw[arrow] (split -| h.north) -- (h.north);
\draw[arrow] (split -| r.north) -- (r.north);

\coordinate (merge) at ($(n7.north)+(0,0.90)$);
\draw[line] (e.south) -- (e.south |- merge);
\draw[line] (h.south) -- (h.south |- merge);
\draw[line] (r.south) -- (r.south |- merge);
\draw[line] (e.south |- merge) -- (r.south |- merge);
\draw[arrow] (merge) -- (n7.north);

\draw[arrow] (n7) -- (n8);
\draw[arrow] (n8) -- (n9);
\end{tikzpicture}
\end{center}

\vspace{0.05cm}
{\cuerpoSmall
Se compararon los métodos de \textbf{Euler}, \textbf{Heun} y
\textbf{Runge--Kutta de cuarto orden (RK4)}. La precisión se evaluó mediante
el error cuadrático medio (RMSE) frente a los datos digitalizados.}

\end{minipage}%
\hfill
% ============================================================================
% COLUMNA 2
% ============================================================================
\begin{minipage}[t]{37.44cm}
\vspace{0pt}
\cuerpo
\raggedright

\bloque{Sistema de ecuaciones diferenciales}

El estado $x(t)\in[0,1]$ cuantifica la adherencia del flujo:
$x=1$ corresponde al extremo adherido y $x=0$ al desprendido.
\begin{align*}
  b\dot{x} &= F(x,\alpha^*), &
  \alpha^* &= \alpha-\tau_2\dot{\alpha},\\[0.18cm]
  F(x,a) &= [x_0(a)-x]\,F_{hm}(x,a)+F_{sd}(x,a), &
  x_0(a) &= \frac{1}{1+(a/\alpha_c)^8}.
\end{align*}

$\alpha^*$ es el ángulo de ataque efectivo, $\tau_2$ el retardo cinemático y
$b>0$ controla la velocidad de evolución del estado. $F_{hm}$ describe la
relajación y $F_{sd}$ una perturbación suave.

\vspace{0.12cm}
El coeficiente de fuerza normal se calcula como
\begin{equation*}
  C_N=C_{N,\mathrm{att}}(\alpha)\,g(x)
     +C_{N,\mathrm{det}}(\alpha)\,[1-g(x)],
  \qquad
  g(x)=\frac{2\sqrt{x}+x}{3},
\end{equation*}
y el movimiento oscilatorio del ala se prescribe mediante
\begin{equation*}
  \alpha(t)=\alpha_m+A\sin(\omega t).
\end{equation*}

{\cuerpoSmall\color{utecGray}
\textbf{Nota:} la formulación integrada en los resultados es la no lineal; la
ecuación lineal del informe se incluye solo como referencia.}

\bloque{Parámetros del modelo}

\begin{center}
\begin{tabular*}{\linewidth}{@{\extracolsep{\fill}} l r @{}}
  \toprule
  \tabhead{Parámetro} & \tabhead{Valor} \\
  \midrule
  Cuerda, $c$                 & $0.242$ m \\
  Velocidad, $U_\infty$       & $40$ m/s \\
  Frecuencia, $\omega$        & $1.2566$ rad/s \\
  Ángulo medio, $\alpha_m$    & $25^\circ$ \\
  Amplitud, $A$               & $24^\circ$ \\
  Ángulo crítico, $\alpha_c$  & $18.657^\circ$ \\
  Constante de relajación, $b$& $0.0422$ s \\
  Retardo cinemático, $\tau_2$& $0.0424$ s \\
  Centro elipse, $x_c$        & $0.3527$ \\
  Semialtura, $x_h$           & $0.3436$ \\
  Inclinación, $\gamma$       & $0.6042$ \\
  Perturbación, $m_x$         & $0.0010$ \\
  \bottomrule
\end{tabular*}
\end{center}

\vspace{0.05cm}
{\cuerpoSmall
$\alpha_w=6.382^\circ$ \textperiodcentered\ $\alpha_{off}=4.378^\circ$
\textperiodcentered\ anchos de $F_{sd}$: $8^\circ$ y $0.5$.
Los valores se obtuvieron por ajuste a los datos digitalizados, valores del
artículo de referencia y supuestos de la implementación.}

\bloque{Resultados principales}

\subbloque{A. Equilibrios estáticos e histéresis}
La continuación numérica identificó \textbf{774 puntos}: 573 estables y 201
inestables. Los pliegues delimitan la zona de histéresis y, para un mismo
$\alpha$, coexisten dos equilibrios estables y uno inestable.

\vspace{0.25cm}
\begin{tikzpicture}
  \fill[utecNavy] (0,0) rectangle (37.44,4.9);
  \foreach \x/\val/\lab in {
    4.68/12.529\textdegree/Reenganche,
    14.04/24.819\textdegree/Pérdida,
    23.40/12.289\textdegree/Ancho de histéresis,
    32.76/0.0358/RMSE estático}
  {
    \node at (\x,3.22) {\headingfont\bfseries\color{white}
      \fontsize{27}{30}\selectfont \val};
    \node at (\x,1.35) {\color{white!82}\fontsize{13.5}{15}\selectfont \lab};
  }
  \foreach \x in {9.36,18.72,28.08}
    \draw[white!30, line width=0.8pt] (\x,0.55) -- (\x,4.35);
\end{tikzpicture}

\subbloque{B. Respuesta dinámica}
Caso base: $\alpha(t)=25^\circ+24^\circ\sin(2\pi\cdot 0.2\,t)$.

\vspace{0.05cm}
\begin{center}
\begin{tabular*}{\linewidth}{@{\extracolsep{\fill}} l r @{}}
  \toprule
  \tabhead{Métrica} & \tabhead{Resultado} \\
  \midrule
  Ángulo de pérdida            & $27.35^\circ$ \\
  Ángulo de reenganche         & $8.78^\circ$ \\
  Separación angular           & $18.57^\circ$ \\
  Área del lazo, $A_L$         & $8.285$ \\
  Máximo de $C_N$              & $1.116$ \\
  \bottomrule
\end{tabular*}
\end{center}

{\cuerpoSmall
El régimen periódico se alcanzó después de tres ciclos de oscilación.}

\subbloque{C. Comparación de métodos numéricos}
\begin{center}
\begin{tabular*}{\linewidth}{@{\extracolsep{\fill}} l r r @{}}
  \toprule
  \tabhead{Método} & \tabhead{Error máximo} & \tabhead{Orden} \\
  \midrule
  Euler & $3.1\times10^{-3}$ & $1$ \\
  Heun  & $8.5\times10^{-6}$ & $2$ \\
  \rowcolor{utecLight}
  \textbf{RK4} & $\mathbf{1.7\times10^{-10}}$ & $\mathbf{4}$ \\
  \bottomrule
\end{tabular*}
\end{center}

{\cuerpoSmall
RK4 presentó el menor error numérico en la comparación realizada.}

\end{minipage}%
\hfill
% ============================================================================
% COLUMNA 3
% ============================================================================
\begin{minipage}[t]{37.44cm}
\vspace{0pt}
\cuerpo
\raggedright

\bloque{Resultados gráficos}

\begin{center}
\begin{minipage}[t]{18.05cm}
  \centering
  \includegraphics[width=\linewidth,height=12cm,keepaspectratio]%
    {Figuras/figura2.png}\\[0.24cm]
  {\captionfont\textbf{Figura 2.} Variable de separación $x$ y coeficiente
   $C_N$: ramas estables, rama inestable y puntos de pliegue.\par}
\end{minipage}%
\hfill
\figimage{Figuras/figura3.png}{3}%
  {Evolución temporal del estado y trayectorias del ciclo base en función
   del ángulo de ataque.}{12cm}

\vspace{0.55cm}

\figimage{Figuras/figura4.png}{4}%
  {Lazos de fuerza normal para tres frecuencias de cabeceo, con amplitud de
   $24^\circ$ y ángulo medio de $25^\circ$.}{12cm}
\hfill
\figimage{Figuras/figura5.png}{5}%
  {Comparación de $C_N$ del modelo con los puntos digitalizados del estudio
   de referencia.}{12cm}
\end{center}

\bloque{Sensibilidad y contraste}

\subbloque{Sensibilidad de los parámetros}
\begin{itemize}[leftmargin=0.78cm, itemsep=0.10cm, topsep=0.05cm,
                  label=\textcolor{utecCyan}{\textbullet}]
  \item Frecuencia de $0.1$ a $0.3$ Hz: el área del lazo aumentó
        aproximadamente un \textbf{53\,\%}.
  \item Duplicar $b$: el área aumentó un \textbf{15.6\,\%}.
  \item Duplicar $\tau_2$: el área aumentó un \textbf{22.6\,\%}.
  \item Semieje de $5.7^\circ$ a $6.7^\circ$: el ancho de histéresis pasó de
        $11.31^\circ$ a $13.27^\circ$.
\end{itemize}

\subbloque{Contraste con datos digitalizados}
\begin{itemize}[leftmargin=0.78cm, itemsep=0.10cm, topsep=0.05cm,
                  label=\textcolor{utecCyan}{\textbullet}]
  \item \textbf{RMSE estático: $0.0358$} (44 puntos).
  \item \textbf{RMSE dinámico: $0.0500$} (31 puntos); $0.0465$ en los
        23 puntos de ajuste y $0.0589$ en los 8 reservados.
\end{itemize}

{\cuerpoSmall\color{utecGray}
Los puntos reservados pertenecen a la misma curva de ajuste: no constituyen
una validación experimental independiente.}

\bloque{Alcance y limitaciones}
{\cuerpoSmall
\begin{itemize}[leftmargin=0.80cm, itemsep=0.12cm, topsep=0.05cm,
                  label=\textcolor{utecCyan}{\textbullet}]
  \item El modelo reduce el campo de flujo a una variable escalar de
        adherencia; no resuelve la estructura espacial del flujo.
  \item Los coeficientes de la función elíptica son fenomenológicos: no
        representan dimensiones físicas del ala.
  \item El RMSE cuantifica el ajuste a datos digitalizados; los puntos
        reservados no constituyen una validación experimental independiente.
\end{itemize}
}

\bloque{Conclusiones}
\begin{itemize}[leftmargin=0.78cm, itemsep=0.10cm, topsep=0.05cm,
                  label=\textcolor{utecCyan}{\textbullet}]
  \item El modelo no lineal representó la histéresis mediante dos ramas
        estables y una inestable; la continuación numérica ubicó los pliegues.
  \item RK4 fue más preciso que Euler y Heun en las pruebas realizadas.
  \item La frecuencia y los parámetros temporales modificaron el área del
        lazo y la respuesta dinámica.
  \item El semieje horizontal de la función elíptica controló principalmente
        el ancho de la histéresis estática.
  \item El modelo reprodujo las tendencias generales; se requieren datos
        adicionales para validar su capacidad predictiva.
\end{itemize}

\bloque{Referencias}
{\fontsize{13.5}{17}\selectfont
\begin{enumerate}[leftmargin=0.62cm, itemsep=0.08cm, topsep=0.02cm]
  \item Goman, M. G., \& Khrabrov, A. N. (1994). State-space representation
        of aerodynamic characteristics of an aircraft at high angles of
        attack. \emph{Journal of Aircraft, 31}(5), 1109--1115.
        \textcolor{utecNavy}{doi:10.2514/3.46618}
  \item Sereez, M., Abramov, N., \& Goman, M. (2024). CFD simulations and
        phenomenological modelling of aerodynamic stall hysteresis of NACA
        0018 wing. \emph{Aerospace, 11}(5), 354.
        \textcolor{utecNavy}{doi:10.3390/aerospace11050354}
  \item Dahlke, J. A., \& Bettinger, R. A. (2024). Practical implementation
        of pseudo-arclength continuation to ensure consistent path direction.
        \emph{Acta Astronautica, 215}, 205--216.
        \textcolor{utecNavy}{doi:10.1016/j.actaastro.2023.12.007}
  \item Ekaterinaris, J. A., \& Platzer, M. F. (1998). Computational
        prediction of airfoil dynamic stall. \emph{Progress in Aerospace
        Sciences, 33}(11--12), 759--846.
        \textcolor{utecNavy}{doi:10.1016/S0376-0421(97)00012-2}
  \item Chen, M., Li, Z., Gao, Z., \& Wang, J. (2024). Influence of boundary
        layer and pressure lag on unsteady aerodynamics of airfoil based on a
        simple semi-empirical dynamic stall engineering model.
        \emph{Wind Energy, 27}(9), 922--935.
        \textcolor{utecNavy}{doi:10.1002/we.2931}
\end{enumerate}
}

\end{minipage}

\end{document}
